\scanpage{382}
\mspage{106}
% archive sheet number: 381
en admettant $\mathrm{Loc}(7\mathrm{F})$), tout en sachant que la structure homotopique droite associée à $W$ est parfaite à tous égards, on n'a à priori aucune maîtrise sur celle-ci, notamment pour le calcul des fibres $W$-homotopiques et la construction \unsure{des|de} carrés $W$-cartésiens.

Même en ramenant toutes ces constructions à $A^{\wedge}$ (\unsure{prenant}, \unsure{tant y étant|tant qu'à faire}, $A=\mathbb{\Delta}$, \deleted{\unsure{)}} i.e. \unsure{se ramenant|en revenant} au contexte semisimplicial), quoiqu'on sache qu'il y a une factorisation \uline{fonctorielle} $f=pi$ avec $\deleted{i\in W_A,\ p\in}\ i\in \mathrm{TC}\subset W_A$, $p\in F\subset\mathrm{Fib}_{W_A}$, je ne vois aucune façon concrète de l'expliciter.
À moins que la considération des chemins infinis, qui (en l'absence de $\mathrm{Loc}(7)$) \unsure{s'avèrent|s'avère} \unsure{inopérants|inopérante} dans $\mathrm{Cat}$, \unsure{sauve|s'avère|s'avérer} \unsure{la mise|là utile} dans $\mathbb{\Delta}^{\wedge}$ ??

